\documentclass[12pt,a4paper]{article}

\usepackage{iftex}
\ifPDFTeX
    \usepackage[utf8]{inputenc}
    \usepackage[T5]{fontenc}
    \usepackage[vietnamese]{babel}
    \usepackage{lmodern}
\else
    \usepackage{fontspec}
    \setmainfont{Times New Roman}
\fi
\usepackage{geometry}
\usepackage{longtable}
\usepackage{booktabs}
\usepackage{array}
\usepackage{enumitem}
\usepackage{hyperref}
\usepackage{xcolor}
\usepackage{amsmath}
\usepackage{amssymb}

\geometry{margin=2.5cm}
\hypersetup{
    colorlinks=true,
    linkcolor=blue,
    urlcolor=blue,
    citecolor=blue
}

\title{\textbf{BÁO CÁO THIẾT KẾ HỆ THỐNG ĐỊNH THỜI LOGIC VÀ XỬ LÝ LUỒNG PHÂN TÁN}}
\author{\textbf{Đặc tả Kiến trúc Hệ thống: Strict Watermark \& Heuristic Watermark}}
\date{}

\begin{document}

\maketitle

\begin{abstract}
Tài liệu này trình bày chi tiết thiết kế học thuật và kiến trúc triển khai thực tế của hệ thống xử lý dòng dữ liệu log phân tán hoạt động liên tục 24/7/365. Hệ thống giải quyết bài toán định thời logic trên event-time trong điều kiện dòng dữ liệu vô hạn, bất tuần tự nghiêm trọng và có độ trễ truyền dẫn biến động lớn. Chúng tôi đề xuất hai cơ chế định thời logic độc lập nhằm tối ưu hóa các mục tiêu đánh đổi khác nhau: (1) Cơ chế \textbf{Strict Watermark} cam kết bảo toàn dữ liệu tuyệt đối (0\% loss) và đảm bảo ngữ nghĩa xử lý đúng một lần (Exactly-Once) phục vụ các nghiệp vụ đối soát tài chính; (2) Cơ chế \textbf{Heuristic Watermark} ứng dụng bộ ước lượng phân vị thích ứng dựa trên cấu trúc dữ liệu \textit{DDSketch} kết hợp đường ống đền bù dữ liệu muộn (DLQ Correction) nhằm đáp ứng các yêu cầu về độ trễ cực thấp (Low Latency). Báo cáo phân tích sâu sắc các khía cạnh về quản lý trạng thái phân tầng (Tiered Storage), cơ chế chịu lỗi phân tán (Fault Tolerance), giám sát chất lượng và quy trình vận hành hệ thống.
\end{abstract}

\newpage
\tableofcontents
\newpage

\section{Tổng quan và Đặt tả bài toán}

\subsection{Bài toán xử lý luồng log phân tán}
Hệ thống xử lý luồng sự kiện (log stream) thu nhận dữ liệu từ các máy chủ Web Server hoạt động trong môi trường phân tán 24/7/365. Dòng dữ liệu này mang ba đặc tính cơ bản:
\begin{itemize}
    \item \textbf{Vô hạn (Unbounded)}: Luồng dữ liệu sinh ra liên tục theo thời gian thực và không có điểm dừng logic.
    \item \textbf{Bất tuần tự nghiêm trọng (Heavy-Tail Out-of-Order)}: Do các sự cố về mạng, kỹ thuật đệm dữ liệu (buffering) tại nguồn phát, hoặc truyền lại TCP, các bản ghi sự kiện có thể đến hệ thống tiêu thụ không theo đúng thứ tự thời gian sinh ra.
    \item \textbf{Tích lũy trạng thái (Stateful)}: Hệ thống yêu cầu phải gom cụm và tính toán thống kê các bản ghi sự kiện theo các cửa sổ thời gian sự kiện (event-time windows) trước khi đưa ra kết quả cuối cùng.
\end{itemize}

\subsection{Cam kết kỹ thuật và Đánh đổi kiến trúc}
Hệ thống được thiết kế để cung cấp khả năng lựa chọn linh hoạt giữa hai hướng thiết kế chính tùy thuộc vào mục tiêu nghiệp vụ:

\begin{longtable}{p{0.3\textwidth} p{0.32\textwidth} p{0.32\textwidth}}
\toprule
\textbf{Tiêu chí cam kết} & \textbf{Chế độ Strict Watermark} & \textbf{Chế độ Heuristic Watermark} \\
\midrule
\endhead
\textbf{Tỷ lệ mất dữ liệu} & Cam kết tuyệt đối 0\% mất dữ liệu. & Kiểm soát trong khoảng $\le 1\%$ (Bursts $\le 5\%$). \\
\textbf{Ngữ nghĩa xử lý} & Exactly-Once (đầu vào và đầu ra). & At-least-once tại lõi, Exactly-once qua Correction. \\
\textbf{Độ trễ xử lý (Latency)} & Trung bình $\approx 15$ giây ($\delta_{base} + W$). & Thấp, trung bình $\approx 5$ giây ($W + L_{eff}$). \\
\textbf{Sự phụ thuộc Upstream} & Yêu cầu Punctuation Token đồng bộ. & Độc lập, không phụ thuộc nguồn phát. \\
\textbf{Độ khả dụng cao (HA)} & Điều phối Raft HA (3 nodes). & Đồng bộ Active-Standby (2 nodes). \\
\textbf{Ứng dụng tiêu biểu} & Đối soát tài chính, kế toán, hóa đơn. & Live Dashboard, hệ thống cảnh báo, Machine Learning. \\
\bottomrule
\caption{Bảng so sánh các đặc tính thiết kế giữa hai chế độ định thời logic.}
\end{longtable}

---

\section{Mô hình thời gian và Phân chia cửa sổ}

\subsection{Trục thời gian logic và Cửa sổ thời gian}
Hệ thống sử dụng thời gian sự kiện (Event-Time), ký hiệu là $T_{event}$, làm trục thời gian logic cố định gắn liền với thời điểm bản ghi được tạo ra tại nguồn phát. Thời gian xử lý vật lý (Processing-Time) tại các Worker hoàn toàn bị loại bỏ khỏi logic phân chia cửa sổ dữ liệu.

Luồng dữ liệu vô hạn được phân chia thành các cửa sổ thời gian cố định dạng Tumbling Window có kích thước $W = 5$ giây. Mỗi bản ghi có thời gian sự kiện $T_{event}$ sẽ được ánh xạ duy nhất vào một cửa sổ $[window\_start, window\_end)$ theo công thức:
\begin{equation}
window\_start = \lfloor T_{event} / W \rfloor \times W
\end{equation}
\begin{equation}
window\_end = window\_start + W
\end{equation}

\subsection{Hàng đợi ưu tiên giới hạn tại Worker (Client-Side Bounded Priority Queue)}
Do cấu trúc phân vùng dữ liệu đầu vào (Kafka Partition) chỉ đảm bảo thứ tự ghi tuần tự theo Offset mà không đảm bảo thứ tự thời gian sự kiện, Worker Node triển khai một hàng đợi ưu tiên giới hạn bằng Min-Heap hoạt động trên bộ nhớ RAM nhằm sắp xếp lại dòng dữ liệu đầu vào theo $T_{event}$:
\begin{itemize}
    \item \textbf{Cơ cấu sắp xếp}: Min-Heap lưu trữ các phần tử theo bộ ba giá trị $(T_{event}, \text{sequence}, \text{event})$ đảm bảo tính ổn định của thứ tự nạp khi các bản ghi trùng timestamp vật lý.
    \item \textbf{Giới hạn dung lượng}: Hàng đợi có sức chứa tối đa là $10.000$ bản ghi nhằm bảo vệ bộ nhớ RAM không bị tràn (OOM). Khi vượt quá giới hạn này, phần tử ở đỉnh heap (có $T_{event}$ nhỏ nhất) sẽ bị ép buộc đẩy (pop) ra để xử lý ngay lập tức.
    \item \textbf{Giới hạn thời gian chờ}: Nhằm ngăn chặn hiện tượng kẹt vô hạn khi dòng dữ liệu đầu vào bị chậm hoặc ngắt quãng, hệ thống quy định mốc thời gian chờ tối đa là $1000$ ms. Nếu khoảng thời gian từ lần đẩy dữ liệu cuối cùng vượt quá mốc này, hệ thống sẽ tiến hành giải phóng bản ghi ở đỉnh heap kể cả khi hàng đợi chưa đầy.
\end{itemize}

---

\section{Thiết kế Kiến trúc Strict Watermark (Bảo toàn dữ liệu 0\%)}

\subsection{Cơ chế Upstream Heartbeat Punctuation}
Để đạt được cam kết 0\% mất dữ liệu mà không chịu ảnh hưởng bởi hiện tượng phân vùng dữ liệu không có hoạt động (Idle Partitions), hệ thống triển khai giao thức Punctuation Token từ phía nguồn phát:
\begin{itemize}
    \item \textbf{Punctuation Token}: Là các thông điệp kiểm soát đặc biệt xen kẽ trong luồng dữ liệu, mang giá trị thời gian cam kết $T_{commit}$. Khi nhận được token này, nguồn phát khẳng định ``tất cả các dữ liệu có mốc thời gian $T_{event} \le T_{commit}$ đã được gửi đi''.
    \item \textbf{Empty Punctuation}: Khi một phân vùng dữ liệu không phát sinh log mới (trạng thái Idle), bộ phát dữ liệu đầu nguồn vẫn phải định kỳ mỗi 1000ms gửi đi một Empty Punctuation Token có giá trị $T_{commit}$ tăng dần theo đồng hồ vật lý nhằm thúc đẩy mốc thời gian logic của hệ thống.
\end{itemize}

\subsection{Bảo đảm đơn điệu và Giám sát Clock Skew}
Tại Worker Node, mốc Watermark cục bộ của phân vùng $P_k$, ký hiệu là $LW_i(P_k)$, được cập nhật trực tiếp theo giá trị $T_{commit}$ của Punctuation Token nhận được. Nhằm đối phó với hiện tượng lệch đồng hồ vật lý (Clock Skew) giữa các nguồn phát dữ liệu đầu nguồn làm $T_{commit}$ bị lùi ngược thời gian, hệ thống áp dụng hai mức bảo vệ:
\begin{enumerate}
    \item \textbf{Worker-side Monotonic Enforcement}: Worker chỉ cập nhật $LW_i(P_k)$ nếu $T_{commit} > last\_T\_commit\_seen$. Ngược lại, thông điệp kiểm soát bị loại bỏ hoàn toàn và ghi nhận một cảnh báo không đơn điệu.
    \item \textbf{Coordinator-side Clock Skew Monitor}: Cụm điều phối theo dõi độ lệch đồng hồ thông qua nhịp tim của Ingestor. Nếu độ lệch vượt ngưỡng 2 giây, hệ thống kích hoạt cảnh báo mức Critical để yêu cầu người vận hành kiểm tra đồng bộ NTP.
\end{enumerate}

\subsection{Điều phối phân tán HA dựa trên Raft (Coordinator HA)}
Hệ thống sử dụng một cụm điều phối gồm 3 instance hoạt động theo cơ chế đồng thuận Raft để duy trì siêu dữ liệu toàn cục (Global Metadata) và loại bỏ điểm lỗi duy nhất (SPOF).

\subsubsection{Thuật toán tính toán Global Watermark}
Định kỳ mỗi 200ms, Coordinator Leader thu thập nhịp tim chứa Watermark cục bộ $LW_i(P_k)$ của toàn bộ các phân vùng đang hoạt động từ các Worker Node. Mốc Watermark toàn cục $W_{global}$ được tính toán bằng giá trị nhỏ nhất của tất cả các Watermark cục bộ:
\begin{equation}
W_{global}(t) = \min_{\forall P_k \in \text{Active}} LW_i(P_k)
\end{equation}
Trong đó, các phân vùng ở trạng thái cảnh báo chậm trễ (STALE) hoặc không hoạt động (IDLE) vẫn \textbf{bắt buộc} phải tham gia vào hàm $\min()$. Chỉ những phân vùng được xác định đã sập hoàn toàn (FAILED) mới bị loại bỏ để đảm bảo tính an toàn toán học chống mất mát dữ liệu.

\subsubsection{Cơ chế Fencing Token bảo vệ hệ thống}
Để chống lại hiện tượng phân mảnh mạng dẫn đến đa Leader (Split-Brain), mọi chỉ thị điều phối từ Coordinator gửi xuống Worker đều mang theo giá trị nhiệm kỳ (fencing term). Worker Node duy trì biến ghi nhận nhiệm kỳ cao nhất từng biết $known\_term$. Khi nhận được chỉ thị có $term < known\_term$, Worker lập tức từ chối thực hiện vì chỉ thị này đến từ một Leader lỗi thời.

\subsection{Quy trình bàn giao phân vùng an toàn (Strict Failback Protocol)}
Khi một Worker Node gặp sự cố và hồi phục lại, nó không được phép tự động tiêu thụ lại dữ liệu từ phân vùng cũ. Quá trình bàn giao quyền sở hữu phân vùng (handoff) được quản lý chặt chẽ thông qua Máy trạng thái phân vùng (Partition State Machine) lưu trữ trong Raft log của cụm điều phối với các trạng thái logic: \texttt{ASSIGNED}, \texttt{REASSIGNING}, \texttt{PAUSED}, và \texttt{ORPHANED}.

\begin{enumerate}
    \item **Yêu cầu bàn giao**: Worker phục hồi gửi tín hiệu đăng ký lên Coordinator.
    \item **Khóa phân vùng**: Coordinator ghi nhận trạng thái \texttt{REASSIGNING} vào Raft log.
    \item **Tạm dừng và chốt chặn**: Coordinator gửi lệnh PAUSE kèm fencing term yêu cầu Worker đang gánh hộ phân vùng dừng tiêu thụ dữ liệu, thực hiện flush toàn bộ trạng thái xuống bộ lưu trữ và xác nhận mốc Offset cuối cùng.
    \item **Chuyển giao quyền lực**: Coordinator cập nhật phân bổ nhóm tiêu thụ Kafka, ghi nhận trạng thái \texttt{PAUSED} vào Raft log.
    \item **Kích hoạt Worker đích**: Coordinator gửi lệnh cho phép Worker hồi phục bắt đầu nạp trạng thái từ checkpoint và thực hiện tiêu thụ từ mốc $\text{Offset} + 1$. Sau khi hoàn tất, trạng thái chuyển về \texttt{ASSIGNED}.
\end{enumerate}

---

\section{Thiết kế Kiến trúc Heuristic Watermark (Độ trễ cực thấp)}

\subsection{Bộ ước lượng phân vị thích ứng DDSketch (Estimator Core)}
Trong kịch bản hệ thống đầu nguồn là các ứng dụng cũ không thể sửa đổi để phát sinh Punctuation Token, hoặc khi yêu cầu về độ trễ cực thấp được đặt lên hàng đầu, cơ chế Heuristic Watermark sử dụng thuật toán DDSketch (Distributable Quantile Sketch) để ước lượng phân phối thực nghiệm của trễ truyền dẫn luồng.

\subsubsection{Cơ chế Log-Scale Buckets và Bảo đảm Sai số}
Với tham số sai số tương đối $\alpha = 0.01$ (1\%), mỗi mẫu trễ $x$ được ánh xạ vào một chỉ số thùng chứa (bucket index) theo thang logarit:
\begin{equation}
\text{bucket}(x) = \lceil \log_\gamma(x) \rceil, \quad \gamma = \frac{1+\alpha}{1-\alpha} = \frac{1.01}{0.99} \approx 1.0202
\end{equation}
Ưu điểm của cấu trúc này là đảm bảo sai số tương đối luôn nằm trong biên giới hạn:
\begin{equation}
\frac{|\hat{q} - q|}{q} \le \alpha
\end{equation}
Đặc tính này cực kỳ phù hợp với phân phối trễ truyền dẫn luồng dữ liệu mang đặc tính đuôi dày (Heavy-Tail), nơi sai số tuyệt đối của các thuật toán truyền thống (như T-digest) sẽ bị phóng đại ở vùng phân vị cao.

\subsubsection{Cơ chế giới hạn bộ nhớ (Hard Cap và Collapse)}
Để duy trì dung lượng bộ nhớ cố định cho mỗi phân vùng $\le 100$ KB, DDSketch áp dụng ba chiến lược:
\begin{itemize}
    \item \textbf{Cắt tỉa đuôi (Tail Truncation)}: Mọi mẫu dữ liệu có trễ vượt quá $MAX\_LAG\_ACCEPTED = 1$ giờ sẽ bị loại khỏi bộ ước lượng và đẩy thẳng vào DLQ.
    \item \textbf{Thu hẹp thùng chứa (Bucket Collapse)}: Khi số lượng thùng chứa vượt quá giới hạn $max\_buckets = 1024$, thuật toán tiến hành gộp các thùng chứa ở vùng giá trị thấp lại với nhau để ưu tiên duy trì độ chính xác ở phân vị cao phục vụ tính toán Watermark.
    \item \textbf{Cấu trúc cửa sổ trượt (Sliding Window Sketch)}: Bộ ước lượng sử dụng cấu trúc gồm 60 sketch con tương ứng với 60 giây lịch sử. Mỗi giây trôi qua, sketch con cũ nhất bị loại bỏ và sketch con mới được khởi tạo, đảm bảo bộ ước lượng luôn bám sát sự thay đổi tức thời của hệ thống.
\end{itemize}

\subsection{Cơ chế Tự thích ứng phân vị (Adaptive Percentile)}
Để đảm bảo an toàn trước các biến động đột ngột về lưu lượng mạng (burst), hệ thống tự động điều chỉnh phân vị mục tiêu $p$:
\begin{equation}
L_{eff} = \text{DDSketch.quantile}(p)
\end{equation}
Bình thường, hệ thống duy trì $p = p_{normal} = 0.99$. Khi phát hiện giá trị phân vị tăng vọt vượt quá $2.0$ lần giá trị trung bình của phút trước đó, hệ thống tự động nâng lên $p = p_{safe} = 0.999$ để mở rộng biên bảo vệ. Hệ thống sẽ tự động hạ lại về mốc cũ sau 5 phút vận hành ổn định.

\subsection{Chiến lược khởi động ấm thích ứng (Cold Start Strategy)}
Khi Worker mới khởi động, bộ ước lượng chưa tích lũy đủ dữ liệu lịch sử để cho ra kết quả thống kê tin cậy. Hệ thống áp dụng chu kỳ khởi động ấm qua các pha:
\begin{itemize}
    \item \textbf{Pha 0 (PHASE\_0)}: Kích hoạt khi thời gian chạy $< 5.0$ giây và số mẫu $< 100$. Worker chỉ thực hiện nạp dữ liệu và hoàn toàn không phát sinh Watermark.
    \item \textbf{Pha 1 (PHASE\_1)}: Kích hoạt khi thời gian chạy $< 10.0$ giây hoặc số mẫu $< 1000$. Watermark được tính toán dựa trên cận an toàn (Conservative Prior):
    \begin{equation}
    L_{eff} = \max(L_{sketch\_quantile}, L_{max\_design\_cap})
    \end{equation}
    Với giá trị cận mặc định $L_{max\_design\_cap} = 60.0$ giây (hoặc nạp từ giá trị lưu trữ lịch sử cuối cùng trước khi tắt máy \texttt{baseline\_history}).
    \item \textbf{Pha thường (NORMAL)}: Khi thời gian chạy $\ge 10.0$ giây VÀ số mẫu $\ge 1000$. Hệ thống sử dụng trực tiếp kết quả phân vị từ DDSketch.
\end{itemize}

\subsection{Đường ống xử lý dữ liệu muộn ngoại vi và Cơ chế đền bù (DLQ Correction)}
Do đặc tính ước lượng heuristic, hệ thống chấp nhận một tỷ lệ nhỏ dữ liệu đến muộn hơn mốc chốt sổ toàn cục $W_{global\_h}$. Mọi dữ liệu muộn này bắt buộc phải được xử lý thông qua đường ống ngoại vi:
\begin{enumerate}
    \item \textbf{Phân loại}: Bản ghi sự kiện đến muộn bị loại khỏi dòng xử lý chính và đẩy vào Kafka topic \texttt{late\_logs\_dlq}.
    \item **Tích lũy lô**: Một tiến trình tiêu thụ dữ liệu muộn (DLQ Consumer) định kỳ đọc các bản ghi này theo lô (hourly hoặc real-time tùy cấu hình).
    \item **Tính toán đền bù**: Tiến trình truy vấn kết quả cửa sổ thời gian đã chốt trước đó từ RocksDB, cộng dồn dữ liệu muộn để tính toán giá trị mới, và sinh ra một thông điệp sửa lỗi (Correction Message) chứa thông tin chênh lệch (Delta).
    \item **Đồng bộ hóa Downstream**: Downstream Consumer tiếp nhận thông điệp sửa lỗi và áp dụng cập nhật theo 1 trong 3 mô hình:
        \begin{itemize}
            \item \textit{Incremental Update}: Cộng trực tiếp giá trị Delta vào bản ghi kết quả cũ.
            \item \textit{Replace}: Ghi đè hoàn toàn kết quả cũ bằng kết quả mới.
            \item \textit{Append with Versioning}: Lưu trữ bản ghi kết quả mới như một phiên bản nâng cấp có chỉ số phiên bản tăng dần ($version++$).
        \end{itemize}
\end{enumerate}

\subsection{Bảo vệ bộ ước lượng trong chế độ chạy lại (Replay-Mode Protection)}
Khi Worker Node thực hiện nạp lại dữ liệu lịch sử từ Kafka (Replay-Mode), giá trị trễ đo được (lag) sẽ tăng cao bất thường so với thời gian xử lý vật lý thực tế. Nếu đưa các mẫu này vào bộ ước lượng, trạng thái thống kê sẽ bị ``nhiễm bẩn'' nghiêm trọng.

\subsubsection{Cơ chế Snapshot và Rollback bộ ước lượng}
Hệ thống thiết lập một Snapshot Manager để quản lý lịch sử trạng thái của bộ ước lượng:
\begin{itemize}
    \item Hệ thống duy trì một hàng đợi vòng tròn chứa tối đa 6 ảnh chụp trạng thái (snapshot) của DDSketch, tương ứng với 60 giây lịch sử gần nhất.
    \item Định kỳ mỗi 10 giây (đồng bộ với chu kỳ checkpoint), hệ thống thực hiện sao chụp toàn bộ cấu trúc thùng chứa của DDSketch và đưa vào hàng đợi.
\end{itemize}

\subsubsection{Thuật toán kiểm soát chế độ Replay}
\begin{itemize}
    \item \textbf{Kích hoạt}: Khi xuất hiện một mẫu dữ liệu có độ trễ lớn hơn $10.0$ lần độ trễ cơ sở trung bình, hệ thống xác định trạng thái Replay. Ngay lập tức, bộ ước lượng thực hiện rollback trạng thái về snapshot an toàn tại thời điểm trước sự cố (ví dụ: mốc thời gian hiện tại trừ đi khoảng đệm 5 giây). Đồng thời, tạm dừng hoàn toàn việc cập nhật các mẫu dữ liệu mới vào DDSketch.
    \item \textbf{Thoát chế độ}: Chế độ Replay chỉ được giải phóng khi độ trễ đo được giảm xuống dưới $2.0$ lần độ trễ cơ sở trung bình và duy trì ổn định liên tục trong 10 giây. Khi thoát, hệ thống cho phép tiếp tục ghi nhận mẫu mới và thực hiện hiệu chuẩn lại độ trễ cơ sở.
\end{itemize}

---

\section{Quản lý Trạng thái và Lưu trữ phân tầng (Tiered Storage)}

Để tối ưu hóa chi phí phần cứng và đảm bảo khả năng phục hồi nhanh chóng, hệ thống triển khai kiến trúc lưu trữ trạng thái phân tầng:

\begin{figure}[h]
\centering
\begin{tabular}{p{0.25\textwidth} p{0.32\textwidth} p{0.33\textwidth}}
\toprule
\textbf{Tầng lưu trữ} & \textbf{Thành phần phần cứng} & \textbf{Chức năng hoạt động} \\
\midrule
\textbf{Tier 1: Hot State} & Bộ nhớ RAM \& SSD NVMe cục bộ của Worker. & Lưu trữ bảng MemTable và các file SST của RocksDB cục bộ. Latency xử lý $< 1$ms. \\
\textbf{Tier 2: Warm State} & Docker Shared Volume (Network SSD). & Lưu trữ các incremental checkpoint và metadata để phục vụ failover nhanh ($< 2$s). \\
\textbf{Tier 3: Cold State} & Object Storage (MinIO / AWS S3). & Lưu trữ lưu trữ lịch sử các cửa sổ đã chốt và các bản sao lưu phục hồi thảm họa. \\
\bottomrule
\end{tabular}
\caption{Kiến trúc lưu trữ trạng thái 3 tầng.}
\end{figure}

\subsection{Cô lập RocksDB Instance theo phân vùng}
RocksDB sử dụng khóa độc quyền trên thư mục dữ liệu qua file \texttt{LOCK}. Nhằm loại bỏ hoàn toàn lỗi tranh chấp khóa dẫn đến crash hệ thống khi cấu hình chạy nhiều Worker container chung một phân vùng chia sẻ, hệ thống bắt buộc cấu hình đường dẫn thư mục RocksDB độc lập cho từng phân vùng vật lý:
$$\text{Path} = \text{/data/rocksdb/partition\_}k\text{/}$$
Khi Coordinator chuyển giao phân vùng $P_k$ từ Node A sang Node B, Node A thực hiện đóng instance RocksDB tương ứng một cách an toàn để nhả khóa, cho phép Node B mở instance này một cách độc quyền mà không xảy ra xung đột.

\subsection{Giao thức giải phóng dữ liệu phân tầng (4-State Eviction Protocol)}
Để đảm bảo tính nhất quán của dữ liệu khi Worker Node gặp sự cố trong quá trình đẩy dữ liệu từ SSD cục bộ lên Object Storage, hệ thống quản lý vòng đời trạng thái của từng cửa sổ đã chốt sổ thông qua máy trạng thái 4 bước:

\begin{figure}[h]
\centering
\texttt{CLOSED} $\rightarrow$ \texttt{UPLOADING} $\rightarrow$ \texttt{UPLOADED} $\rightarrow$ \texttt{PURGED}
\caption{Các bước chuyển đổi trạng thái của cửa sổ trong quy trình Eviction.}
\end{figure}

\begin{itemize}
    \item \textbf{CLOSED}: Cửa sổ đã kết thúc tính toán thống kê, kết quả được lưu trữ cục bộ tại Tier 1.
    \item \textbf{UPLOADING}: Tiền trình chạy nền bắt đầu tải kết quả lên MinIO. Trạng thái này được ghi nhận kèm mã giao dịch duy nhất.
    \item \textbf{UPLOADED}: Nhận phản hồi thành công (HTTP 200) từ MinIO. Trạng thái này lưu kèm mã ETag của đối tượng lưu trữ để làm bằng chứng đối soát.
    \item \textbf{PURGED}: Hệ thống thực hiện xóa sạch dữ liệu thô của cửa sổ khỏi RocksDB cục bộ tại Tier 1 để giải phóng không gian đĩa NVMe.
\end{itemize}
Trong quá trình phục hồi sau sự cố, nếu hệ thống phát hiện một cửa sổ ở trạng thái lấp lửng \texttt{UPLOADING}, nó sẽ thực hiện truy vấn khóa lưu trữ trên MinIO bằng mã giao dịch duy nhất. Nếu đối tượng đã tồn tại, trạng thái được tự động chuyển sang \texttt{UPLOADED} để thực hiện xóa local đĩa; ngược lại, hệ thống sẽ thực hiện tải lên lại từ đầu.

---

\section{Hoạch định Năng lực và Triển khai Vận hành}

\subsection{Đánh đổi Kiến trúc theo Định lý CAP/PACELC và Bảo đảm SLA}
Kiến trúc hệ thống xử lý luồng dữ liệu phân tán này được xây dựng dựa trên sự kết hợp chặt chẽ giữa lý thuyết hệ thống phân tán và yêu cầu thực tế của hạ tầng. Dưới đây là phân tích học thuật chi tiết dựa trên Định lý CAP và mô hình mở rộng PACELC (Partition, Availability, Consistency, Else, Latency, Consistency) để làm rõ các đánh đổi thiết kế và cam kết SLA của hệ thống.

\subsubsection{Mô hình hóa PACELC cho Tuyến Strict và Heuristic}
Định lý PACELC chỉ ra rằng: Nếu xảy ra phân mảnh mạng (\textit{Partition} - P), hệ thống phải đánh đổi giữa tính khả dụng (\textit{Availability} - A) và tính nhất quán (\textit{Consistency} - C). Ngược lại (\textit{Else} - E), khi hệ thống vận hành bình thường không có lỗi mạng, hệ thống phải đánh đổi giữa độ trễ (\textit{Latency} - L) và tính nhất quán (\textit{Consistency} - C).

Hệ thống của chúng tôi tích hợp song song hai tuyến xử lý với các cấu hình đánh đổi bổ trợ nhau hoàn hảo:
\begin{enumerate}
    \item \textbf{Tuyến Strict Path (Phân loại PC/EC - Partition Consistency / Else Consistency)}:
    \begin{itemize}
        \item \textbf{Khi xảy ra phân mảnh mạng (P)}: Tuyến này hoạt động theo mô hình \textbf{CP}. Sự an toàn dữ liệu được đặt lên hàng đầu. Khi một Worker Node bị ngắt kết nối hoặc mất liên lạc với Leader Coordinator (thiếu Quorum Raft), toàn bộ tiến trình xử lý trên các phân vùng Kafka tương ứng sẽ lập tức bị đóng băng. Hệ thống từ chối sinh thêm Watermark ảo để tránh ghi đè dữ liệu sai lệch. Việc này chấp nhận dừng hệ thống (hy sinh tính khả dụng $A$) để đảm bảo tính nhất quán dữ liệu ($C$) là tuyệt đối.
        \item \textbf{Khi hoạt động bình thường (E)}: Tuyến này hoạt động theo mô hình \textbf{EC}. Để bảo toàn thứ tự Exactly-Once và 0\% mất dữ liệu, hệ thống bắt buộc phải áp dụng thuật toán Min-of-Min trên toàn bộ phân vùng. Điều này có nghĩa là tiến độ Watermark toàn cục bị kéo chậm bởi phân vùng chậm nhất (Straggler). Hệ thống chấp nhận tăng độ trễ xử lý ($L$) để đạt được sự đồng nhất và chính xác hoàn hảo về mặt trạng thái dữ liệu ($C$).
    \end{itemize}
    
    \item \textbf{Tuyến Heuristic Path (Phân loại PA/EL - Partition Availability / Else Latency)}:
    \begin{itemize}
        \item \textbf{Khi xảy ra phân mảnh mạng (P)}: Tuyến này hoạt động theo mô hình \textbf{AP}. Khi có sự cố phân mảnh mạng, các Worker Node tự trị tiếp tục hoạt động độc lập và tính toán cục bộ dựa trên dữ liệu hiện có trong bộ ước lượng DDSketch của nó mà không chờ đợi sự đồng thuận toàn cục. Dòng chảy dữ liệu tiếp diễn liên tục (bảo toàn tính khả dụng $A$), chấp nhận việc các cửa sổ kết quả chốt sổ có thể bị lệch pha hoặc thiếu dữ liệu (hy sinh tính nhất quán $C$).
        \item \textbf{Khi hoạt động bình thường (E)}: Tuyến này hoạt động theo mô hình \textbf{EL}. Mục tiêu tối thượng là tối thiểu hóa độ trễ ($L$). Hệ thống sử dụng giá trị ước lượng phân vị trễ $p$ từ DDSketch để chủ động chốt sổ sớm cửa sổ thời gian. Các dữ liệu đến muộn hơn mốc chốt sổ sẽ bị loại bỏ khỏi luồng chính và xử lý bù trừ bất đồng bộ qua DLQ, chấp nhận tính nhất quán tạm thời bị vi phạm tại Downstream (chờ sửa lỗi sau) để đổi lấy độ trễ cực thấp ($L$).
    \end{itemize}
\end{enumerate}

\subsubsection{Bảng so sánh đối chiếu đặc tính lý thuyết}
Dưới đây là bảng tổng hợp các tham số đánh đổi thiết kế giữa hai tuyến xử lý để làm cơ sở đối soát học thuật:

\begin{table}[H]
\centering
\small
\begin{tabular}{l p{0.36\textwidth} p{0.38\textwidth}}
\toprule
\textbf{Đặc tính so sánh} & \textbf{Tuyến Strict Path (PC/EC)} & \textbf{Tuyến Heuristic Path (PA/EL)} \\
\midrule
\textbf{Mô hình CAP} & CP (Consistency / Partition Tolerance) & AP (Availability / Partition Tolerance) \\
\textbf{Mô hình PACELC} & PC/EC (Consistency dưới cả hai trạng thái) & PA/EL (Availability khi P, Latency khi E) \\
\textbf{Độ chính xác dữ liệu} & Cam kết 100\% chính xác, 0\% thất thoát & Chấp nhận thất thoát nhất thời $\le 1\%$, sửa lỗi qua DLQ \\
\textbf{Cơ chế tiến Watermark} & Đồng bộ toàn cục (Min-of-Min) qua Coordinator & Tự trị cục bộ (DDSketch quantile) trên Worker \\
\textbf{Ảnh hưởng bởi Straggler} & Bị nghẽn bởi nút chậm nhất toàn hệ thống & Hoàn toàn miễn nhiễm, tự động cô lập đòn bẩy trễ \\
\textbf{RTO (Recovery Time Objective)} & RTO $\approx 10$ -- 15 giây (phụ thuộc bầu cử Raft) & RTO $\le 2$ giây (Aggregator Active-Standby failover) \\
\bottomrule
\end{tabular}
\caption{Bảng so sánh đối chiếu đặc tính PACELC giữa Strict Path và Heuristic Path.}
\label{tab:pacelc_comparison}
\end{table}

\subsubsection{Đảm bảo SLA chất lượng dịch vụ (SLA Guarantees)}
Dựa trên sự phân tách kiến trúc PACELC, hệ thống cam kết các chỉ số SLA kỹ thuật như sau:
\begin{itemize}[leftmargin=*]
    \item \textbf{SLA Tuyến Strict Path}:
    \begin{itemize}
        \item \textit{Độ chính xác}: Cam kết Exactly-Once processing và $0.0\%$ thất thoát dữ liệu đối với mọi bản ghi được ghi nhận thành công vào Kafka.
        \item \textit{Độ trễ chốt sổ}: Độ trễ $P_{99} \le 15.0$ giây trong điều kiện tải bình thường và độ lệch đồng hồ vật lý giữa các node $\le 100$ms. Trong trường hợp có phân mảnh mạng, độ trễ có thể tăng vô hạn cho đến khi khôi phục kết nối.
    \end{itemize}
    \item \textbf{SLA Tuyến Heuristic Path}:
    \begin{itemize}
        \item \textit{Độ trễ chốt sổ}: Cam kết độ trễ chốt sổ $P_{99} \le 5.0$ giây kể cả khi xuất hiện hiện tượng giật lag đường truyền (network spikes) lên tới $2.0$ lần độ trễ trung bình.
        \item \textit{Tính toàn vẹn cuối cùng}: Tỷ lệ dữ liệu bị loại bỏ vào DLQ $\le 1.0\%$ trong trạng thái vận hành ổn định. Thời gian hội tụ và đền bù dữ liệu muộn qua DLQ Correction xuống Downstream đảm bảo $T_{\mathrm{compensation}} \le 5.0$ phút kể từ thời điểm phát hiện dữ liệu muộn.
    \end{itemize}
\end{itemize}

\subsubsection{Đánh đổi Tài nguyên hệ thống}
Bên cạnh đánh đổi về mặt dữ liệu, hệ thống còn thực hiện tối ưu hóa sâu ở mức tài nguyên vật lý:
\begin{itemize}[leftmargin=*]
    \item \textbf{RAM vs. SSD IOPS}: Việc cô lập các RocksDB Instance độc lập cho từng phân vùng Kafka giúp loại bỏ lỗi chiếm dụng file khóa \texttt{LOCK} đồng thời triệt tiêu nguy cơ tràn bộ nhớ RAM (OOM) khi hệ thống chịu áp lực tải cao (Backpressure). Sự đánh đổi là sự gia tăng tải đọc/ghi ngẫu nhiên trên ổ đĩa SSD NVMe (IOPS), được bù đắp bằng việc sử dụng SSD cấp doanh nghiệp có độ bền cao.
    \item \textbf{Băng thông mạng vs. Chi phí lưu trữ}: Mô hình Tiered Storage di chuyển các tệp trạng thái đã chốt sổ từ local SSD (Tier 1 và 2) lên MinIO (Tier 3), giúp hệ thống hoạt động ổn định mà không lo hết dung lượng đĩa cục bộ. Sự chuyển dịch này đánh đổi việc gia tăng băng thông mạng nội bộ khi thực hiện giao thức giải phóng dữ liệu 4 bước (\texttt{CLOSED} $\rightarrow$ \texttt{UPLOADING} $\rightarrow$ \texttt{UPLOADED} $\rightarrow$ \texttt{PURGED}).
\end{itemize}


\subsection{Kịch bản xử lý sự cố khẩn cấp (Operations Runbook)}

\subsubsection{Sự cố 1: Xuất hiện tỷ lệ mất dữ liệu lớn hơn 0 trong chế độ Strict}
\begin{itemize}
    \item \textbf{Triệu chứng}: Hệ thống phát tín hiệu cảnh báo Critical đỏ. Downstream nhận thiếu dữ liệu đối soát.
    \item \textbf{Quy trình xử lý}:
        \begin{enumerate}
            \item Kiểm tra chỉ số \texttt{watermark\_lag} để xác định xem watermark toàn cục có đang bị đóng băng hay nhảy vọt bất thường.
            \item Truy vấn số lượng sự kiện bị loại bỏ trên metric \texttt{non\_monotonic\_punctuation} để xác định xem có lỗi lệch đồng hồ vật lý giữa các nguồn phát dữ liệu đầu nguồn hay không.
            \item Kiểm tra nhịp tim của các Ingestor trên giao thức gRPC để xác định xem có nguồn phát nào bị sập mà không kích hoạt Empty Punctuation.
            \item Tiến hành đồng bộ đồng hồ các máy chủ nguồn phát bằng dịch vụ NTP/Chrony, thực hiện khởi động lại các Ingestor bị stuck.
        \end{enumerate}
\end{itemize}

\subsubsection{Sự cố 2: Worker Node bị sập do lỗi tràn bộ nhớ (OOM)}
\begin{itemize}
    \item \textbf{Triệu chứng}: Container của Worker bị kill đột ngột bởi hệ điều hành. Chỉ số sống sót \texttt{worker\_alive\_count} giảm.
    \item \textbf{Quy trình xử lý}:
        \begin{enumerate}
            \item Xác định xem cơ chế chống nghẽn (Backpressure) có hoạt động kịp thời hay không bằng cách đối chiếu metric \texttt{backpressure\_pause\_rate}.
            \item Nếu cơ chế pause chậm trễ do kích thước lô nạp từ Kafka quá lớn, thực hiện cấu hình giảm biến môi trường \texttt{BP\_PAUSE\_THRESHOLD} từ 500 xuống 300 bản ghi.
            \item Thực hiện giảm dung lượng bộ đệm RocksDB Block Cache hoặc cấu hình giới hạn kích thước nạp tối đa từ phía Kafka (\texttt{fetch.max.bytes}).
            \item Kích hoạt tái phân bổ phân vùng để các Worker còn sống tiếp quản trạng thái từ Tier 2 SSD.
        \end{enumerate}
\end{itemize}

\end{document}
